\documentclass[11pt,a4paper]{article}

\usepackage[T1]{fontenc}
\usepackage[utf8]{inputenc}
\usepackage{booktabs}
\usepackage{geometry}
\usepackage{hyperref}
\usepackage{microtype}
\usepackage{array}
\usepackage{xcolor}

\geometry{margin=1in}
\hypersetup{
  colorlinks=true,
  linkcolor=blue,
  urlcolor=blue,
  citecolor=blue
}

\title{Thai Speech Anti-Spoofing with LFCC and a Small CNN}
\author{Biometrics Mini-Project}
\date{April 2026}

\begin{document}

\maketitle

\begin{abstract}
This project studies Thai speech anti-spoofing as a binary classification task. Each input utterance is classified as either genuine speech or spoofed speech. The current experiment uses the local AI For Thai spoofing data under \texttt{data/}, extracts LFCC acoustic features, and trains a lightweight convolutional neural network. The final 1,200-file subset is balanced at the binary-class level and includes three spoof sources: VAJA text-to-speech, F0-40 modification, and F0-160 modification. On the held-out test split, the LFCC + \texttt{small\_cnn} model reached 0.8817 accuracy, 0.8819 balanced accuracy, 0.8842 F1-score, and 0.1236 Equal Error Rate (EER). These results provide a realistic baseline for the expanded Thai anti-spoofing dataset.
\end{abstract}

\section{Introduction}

Voice biometric systems can be attacked using synthetic speech, converted speech, or signal-level speech modifications. Anti-spoofing systems are designed to detect these attacks before a biometric matcher accepts an utterance.

This mini-project builds a practical Thai speech anti-spoofing pipeline. The goal is not to train on the full dataset, but to create a reproducible local experiment that uses real Thai audio folders, configurable paths, cached feature extraction, a lightweight neural model, and standard evaluation metrics.

\section{Dataset}

The experiment uses the AI For Thai ``Spoof Detection for Speech Signal (Thai language)'' data placed under \texttt{data/}. The local dataset contains one genuine speech source and three spoof sources.

\begin{table}[h]
\centering
\caption{Usable local dataset sources after skipping \texttt{\_\_MACOSX} folders and \texttt{.\_*} metadata files.}
\label{tab:dataset-sources}
\begin{tabular}{llr}
\toprule
Folder & Role / attack type & Usable files \\
\midrule
\texttt{Corpus-Spoof-genuine} & Genuine speech & 4,583 \\
\texttt{Corpus-Spoof-VAJA} & VAJA TTS spoof & 4,583 \\
\texttt{Corpus-Spoof-F0\_40} & F0-40 modification spoof & 4,582 \\
\texttt{Corpus-Spoof-F0\_160\_new} & F0-160 modification spoof & 4,583 \\
\midrule
Total & & 18,331 \\
\bottomrule
\end{tabular}
\end{table}

A 1,200-file subset was used for the final run so the experiment would remain practical on local hardware. The spoof side was attack-balanced across VAJA, F0-40, and F0-160. Files sharing the same base utterance ID were kept on the same side of the split. For example, genuine \texttt{thai0001}, VAJA \texttt{thai\_000001}, and F0 variants such as \texttt{f0\_40\_thai0001} must not be split across train and test.

\begin{table}[h]
\centering
\caption{Current manifest split for \texttt{aiforthai\_lfcc\_small\_cnn\_1200}.}
\label{tab:manifest-split}
\begin{tabular}{lrrrr}
\toprule
Split & Genuine & VAJA spoof & F0-40 spoof & F0-160 spoof \\
\midrule
Train feature split & 510 & 168 & 168 & 168 \\
Test feature split & 90 & 32 & 32 & 32 \\
\bottomrule
\end{tabular}
\end{table}

The feature-level training split contains 510 genuine and 504 spoof files. The feature-level test split contains 90 genuine and 96 spoof files. During model training, the training split is divided internally into 811 training samples and 203 validation samples using a 20\% validation fraction.

\section{Method}

The pipeline has four stages:

\begin{enumerate}
  \item Discover audio files and infer binary labels from the folder names.
  \item Build a group-safe train/test manifest with balanced spoof attacks.
  \item Extract LFCC feature matrices and cache them as pickle files.
  \item Train and evaluate a lightweight CNN classifier.
\end{enumerate}

\subsection{Feature Extraction}

The current experiment uses LFCC only. Each audio file is converted to mono audio at 16 kHz. LFCC features are extracted with a 30 ms analysis window, 15 ms hop, 1024-point FFT, 70 linear filters, 20 cepstral coefficients, log energy, delta features, and delta-delta features. Each utterance is converted to a fixed \texttt{128 x 60} feature matrix by repeat-padding short samples and truncating long samples.

\subsection{Model}

The model is \texttt{small\_cnn}, a compact TensorFlow/Keras convolutional neural network. It uses three convolution blocks with batch normalization and max pooling, followed by global average pooling, dropout, a dense hidden layer, and a sigmoid output. Spoof speech is treated as the positive class.

\begin{table}[h]
\centering
\caption{Final experiment configuration.}
\label{tab:experiment-config}
\begin{tabular}{ll}
\toprule
Setting & Value \\
\midrule
Config file & \texttt{ThaiSpoof/configs/aiforthai\_lfcc\_small\_cnn\_1200.json} \\
Run directory & \texttt{ThaiSpoof/runs/aiforthai\_lfcc\_small\_cnn\_1200} \\
Feature & LFCC \\
Model & \texttt{small\_cnn} \\
Input shape & \texttt{128 x 60 x 1} \\
Batch size & 8 \\
Maximum epochs & 30 \\
Learning rate & 0.0001 \\
Early-stopping patience & 6 \\
\bottomrule
\end{tabular}
\end{table}

\section{Evaluation Metrics}

The project reports accuracy, balanced accuracy, precision, recall, F1-score, Equal Error Rate (EER), and confusion matrix counts. Spoof speech is the positive class, so true positives are spoof files correctly classified as spoof.

\section{Results}

The final metrics were read from:

\begin{center}
\texttt{ThaiSpoof/runs/aiforthai\_lfcc\_small\_cnn\_1200/results/lfcc\_small\_cnn\_metrics.csv}
\end{center}

\begin{table}[h]
\centering
\caption{LFCC + \texttt{small\_cnn} metrics.}
\label{tab:metrics}
\begin{tabular}{lrrrrrr}
\toprule
Split & Accuracy & Balanced accuracy & Precision & Recall & F1 & EER \\
\midrule
Train & 0.9778 & 0.9778 & 0.9800 & 0.9752 & 0.9776 & 0.0222 \\
Validation & 0.9015 & 0.9013 & 0.9355 & 0.8614 & 0.8969 & 0.0985 \\
Test & 0.8817 & 0.8819 & 0.8936 & 0.8750 & 0.8842 & 0.1236 \\
\bottomrule
\end{tabular}
\end{table}

\begin{table}[h]
\centering
\caption{Test split confusion matrix. Spoof is the positive class.}
\label{tab:confusion}
\begin{tabular}{lrr}
\toprule
 & Predicted genuine & Predicted spoof \\
\midrule
True genuine & 80 & 10 \\
True spoof & 12 & 84 \\
\bottomrule
\end{tabular}
\end{table}

On the 186-file test split, the model correctly classified 80 of 90 genuine files and 84 of 96 spoof files. The test balanced accuracy was 0.8819, which is close to the standard accuracy because the test split is nearly balanced.

\section{Discussion}

The expanded AI For Thai subset is more informative than the earlier smaller experiments. It includes three spoof sources rather than a single attack family, and the spoof side is balanced across VAJA, F0-40, and F0-160. This makes the task harder and gives a more realistic baseline for Thai anti-spoofing.

The gap between train performance and validation/test performance suggests that the model learns the training subset well, but still has limited generalization to held-out utterance groups. This is expected for a compact CNN trained on a 1,200-file subset. The result should be interpreted as a baseline rather than a final anti-spoofing system.

The confusion matrix shows both error types. Ten genuine files were incorrectly marked as spoof, and twelve spoof files were incorrectly marked as genuine. In a biometric security setting, the false negatives are especially important because they represent spoofed speech that passed through the detector.

\section{Limitations}

\begin{itemize}
  \item The experiment uses a 1,200-file subset, not the full available dataset.
  \item Only LFCC + \texttt{small\_cnn} was evaluated in the current final run.
  \item The test split contains the same attack families as the training split: VAJA, F0-40, and F0-160.
  \item Generalization to unseen attacks was not measured in the current run.
  \item Fixed-size feature matrices truncate long utterances and repeat short utterances.
  \item Some dataset archive metadata files had to be skipped because they are not valid audio files.
\end{itemize}

\section{Reproducibility}

The current experiment can be reproduced with:

\begin{verbatim}
python -m ThaiSpoof.project.run_experiment ^
  --config ThaiSpoof/configs/aiforthai_lfcc_small_cnn_1200.json ^
  --stage all --force-split --force-extract
\end{verbatim}

The main generated artifacts are:

\begin{itemize}
  \item \texttt{manifest.csv}
  \item \texttt{features/lfcc/LFCC\_Train\_genuine\_510.pkl}
  \item \texttt{features/lfcc/LFCC\_Train\_spoof\_504.pkl}
  \item \texttt{features/lfcc/LFCC\_Test\_genuine\_90.pkl}
  \item \texttt{features/lfcc/LFCC\_Test\_spoof\_96.pkl}
  \item \texttt{models/lfcc\_small\_cnn.keras}
  \item \texttt{results/lfcc\_small\_cnn\_history.csv}
  \item \texttt{results/lfcc\_small\_cnn\_metrics.csv}
\end{itemize}

\section{Conclusion}

This project produced a reproducible Thai speech anti-spoofing baseline using AI For Thai audio, LFCC features, and a compact CNN classifier. The final 1,200-file run achieved 0.8817 test accuracy, 0.8842 test F1-score, and 0.1236 EER. Because the dataset includes VAJA, F0-40, and F0-160 spoof sources, the result is a more realistic baseline than earlier single-source experiments.

Future work should train on a larger subset, evaluate each attack family separately, reserve at least one attack type as an unseen-attack test, and compare LFCC against other acoustic or learned front ends.

\end{document}
